\documentclass[11pt]{article}
\usepackage[margin=0.88in]{geometry}
\usepackage[T1]{fontenc}
\usepackage{lmodern}
\usepackage{amsmath,amssymb,amsthm,mathtools}
\usepackage{booktabs,longtable,array,microtype,xcolor}
\usepackage{hyperref}
\hypersetup{colorlinks=true,linkcolor=blue!45!black,urlcolor=blue!55!black,
 pdftitle={A twelve-variable equation for Con(ZFC)},
 pdfauthor={Research report prepared for Vladimir Reshetnikov}}
\setlength{\emergencystretch}{3em}
\setlength{\parskip}{3pt}
\newtheorem{theorem}{Theorem}[section]
\newtheorem{lemma}[theorem]{Lemma}
\newtheorem{proposition}[theorem]{Proposition}
\newtheorem{corollary}[theorem]{Corollary}
\theoremstyle{remark}\newtheorem{remark}[theorem]{Remark}
\newcommand{\Z}{\mathbb Z}
\newcommand{\Q}{\mathbb Q}
\newcommand{\Pos}{\mathbb Z_{>0}}
\newcommand{\Con}{\operatorname{Con}}
\newcommand{\Norm}{\operatorname{N}}
\newcommand{\pc}{\operatorname{pc}}
\newcommand{\topdeg}{\mathrm{top}}
\newcommand{\pin}{\texttt{f0ae4e0991ca44a7bb4e5671d42e33d775e48035}}
\title{\textbf{A twelve-variable equation for \(\Con(\mathrm{ZFC})\)}\\[5pt]
\large Quadratic norms, rational Pell integrality,\\
and a smaller fixed-coefficient compiler}
\author{Research report prepared for Vladimir Reshetnikov\\
\small OpenAI-assisted mathematical research for ProveIt}
\date{22 October 2026}
\begin{document}
\maketitle
\begin{abstract}
We combine and extend the seven current ProveIt constructions of
Diophantine formulations of ZFC consistency. The main result is an
explicit fixed-coefficient polynomial in twelve positive integer
unknowns, of exact degree \(1353\), evaluated by \(116\) binary arithmetic
operations. Its absence of positive solutions is provably equivalent
in ZFC to \(\Con(\mathrm{ZFC})\). The previous explicit construction
in this compiler family used thirteen positive unknowns.
A three-equation presentation has maximum degree \(1344\) and costs
\(112\) operations. A separately justified signed-domain substitution
gives a polynomial in \(26\) unrestricted integers, of degree \(2657\)
and \(168\) operations, improving the current \(29\)-integer construction.
Independently, a stricter local language and denser synchronization
layout reduce the fixed base from \(2^{5^{18}}\) to \(2^{5^{15}}\):
the binary length of \(B-1\) is \(125\) times smaller.
All formulas, coefficient recipes, circuit generators, literal
transition data, proof notes, and verification scripts accompany this
report. The new results are ordinary mathematical proofs, with
independent algebraic and finite checks; they are not new end-to-end
Lean or Rocq formalizations. We claim improvements to this explicit
family, not a global minimum or a proof of consistency.
\end{abstract}
\tableofcontents
\newpage

\section{Results, scope, and comparison}

The source snapshot reviewed here is
\begin{center}\small\pin.\end{center}
The principal source directory is
\begin{center}\small
\texttt{SetTheory/Cardinals/docs/reports/}\\
\texttt{hilbert-tenth-problem/con-zfc-diophantine-forms/}.
\end{center}
At this snapshot it contains packages 09--15, which supersede packages
01--08 \cite{repo}. We use the repository's package numbers as source identifiers,
without assigning priority beyond the repository's own reconciliation.
The comparison concerns actual finite polynomial descriptions, including
exact recipes for their fixed integer coefficients.

\subsection{What the theorem says}

Fix the syntactic proof predicate for ZFC used in the K215 calibration.
Then \(\Con(\mathrm{ZFC})\) means that there is no finite derivation of
falsity from those axioms. Our main polynomial \(P_{12}\) satisfies
\[
 \mathrm{ZFC}\ \vdash\
 \left(\exists \mathbf x\in\Pos^{12}\ P_{12}(\mathbf x)=0\right)
 \ \longleftrightarrow\
 \neg\Con(\mathrm{ZFC}).
\]
Equivalently,
\begin{equation}\label{eq:main-con}
 \boxed{\ \mathrm{ZFC}\ \vdash\
 \left(\forall \mathbf x\in\Pos^{12}\ P_{12}(\mathbf x)\ne0\right)
 \ \longleftrightarrow\ \Con(\mathrm{ZFC})\ .\ }
\end{equation}
The four coefficient ports occurring below are fixed integers, not
additional quantified unknowns. Strict positivity is part of this
statement. A construction over all integers is provided separately in
Section~\ref{sec:integer}.

The proof factors through the exact blank run of K215:
\[
 \text{Diophantine solution}
 \ \longleftrightarrow\
 \text{blank K215 halts}
 \ \longleftrightarrow\
 \text{a finite ZFC refutation exists}.
\]
The second equivalence has existing Lean and Rocq proofs for the literal
machine \cite{endpoint}. The first uses the existing inputless compiler and the new
mathematical arguments supplied here. We do not infer an internal ZFC
derivation merely from a proof assistant's external correctness record:
the ordinary arithmetic, finite-computation, and calibration arguments
are the arguments to be formalized in ZFC.

\subsection{The measured improvements}

An operation is one binary \(+\), \(-\), or \(\times\), including
multiplication by a fixed numeral. Reuse and the comparisons in a system
are free. Degree is total degree in the quantified unknowns. Coefficient
bit length is a separate metric.

\begin{center}\small
\begin{tabular}{p{0.38\linewidth}rrrr}
\toprule
Construction & Unknowns & Equations & Degree & Ops\\
\midrule
Existing package 09 product & \(13\) positive & 1 & 571 & 84\\
Existing package 09 lower-degree product & \(13\) positive & 1 & 515 & 85\\
\textbf{New norm system} & \(\mathbf{12}\) positive & 3 & 1344 & 112\\
\textbf{New norm product} & \(\mathbf{12}\) positive & 1 & 1353 & 116\\
New lower-degree norm product & \(12\) positive & 1 & 1177 & 123\\
New norm sum of squares & \(12\) positive & 1 & 2688 & 119\\
Existing package 14 normalized product & \(29\) integer & 1 & 300 & 119\\
\textbf{New integer product} & \(\mathbf{26}\) integer & 1 & 2657 & 168\\
New integer system & \(26\) integer & 3 & 2640 & 164\\
New integer sum of squares & \(26\) integer & 1 & 5280 & 171\\
\bottomrule
\end{tabular}
\end{center}

These variable reductions cost degree and arithmetic operations.
The existing \(67\)-operation, \(15\)-positive-variable product remains
the best stated operation count in the reviewed family \cite{repo,blank}. The new
coefficient recipe can also be used with that older product, without
changing its degree or operation count.

\begin{center}\small
\begin{tabular}{lrr}
\toprule
Coefficient metric & Existing T4 & New restricted T4\\
\midrule
Local tiles & 37,075 & 5,280\\
Radix exponent & \(5^8\) & \(5^7\)\\
Cell length & \(5^{10}\) & \(5^8\)\\
Base \(B\) & \(2^{5^{18}}\) & \(2^{5^{15}}\)\\
Bit length of \(B-1\) & 3,814,697,265,625 & 30,517,578,125\\
\bottomrule
\end{tabular}
\end{center}

The coefficient descriptions are finite and exact. The final integers
are not expanded: even the new \(B-1\) needs about \(3.81\) decimal GB
in packed binary. An arithmetic circuit's small description therefore
must not be confused with small coefficient height.

\subsection{Which existing ideas are combined}

\begin{center}\small
\begin{tabular}{>{\raggedright\arraybackslash}p{0.16\linewidth}>{\raggedright\arraybackslash}p{0.35\linewidth}>{\raggedright\arraybackslash}p{0.39\linewidth}}
\toprule
Source & Contribution retained & Role in the new result\\
\midrule
09 & Weighted amalgamation of three Pell roots; weak unit-sign recovery &
Starting point for absorbing a fourth coordinate\\
10, 13 & Root elimination and transfer of a coherent coefficient tuple &
Separation of arithmetic from the local compiler\\
11 & Head-filtered four-cell T4 language; signed main-root proof &
Base for the smaller coefficient language and layout\\
12 & Ordinary strong Pell interface; signed soundness; boundary bootstrap &
Soundness of reconstructed roots without assumed signs\\
14 & Seven positive ports, odd multiplier, square power scale &
The \(26\)-integer version\\
15 & Explicit low-degree alternatives and careful circuit comparisons &
Tradeoff benchmark; no claim to improve its low degrees\\
\bottomrule
\end{tabular}
\end{center}

The new step is to encode the first root, main root, strong root, and
auxiliary root through one supplied integer, while retaining the main
root's required congruence. A rational-Pell integrality lemma closes the
denominator gap that a bare resultant would leave.

\section{The explicit twelve-variable polynomial}\label{sec:formula}

Supply exactly the twelve positive integers
\[
 (J,F,\alpha,z,h,i,T,s,w,\eta,\zeta,r)\in\Pos^{12}.
\]
Fix one complete coefficient tuple
\[
 (\beta,K_0,M_C,M_F),\qquad
 \beta=B-1,\quad K_0=D_C+BD_R,\quad
 M_F=M_F^{\rm nat}+B-1.
\]
The new default tuple is defined in Section~\ref{sec:coefficients}.
The original T4 tuple is also valid as a complete alternative. In
particular the fourth port is the \emph{shifted} mask.

Define polynomial abbreviations
\begin{align*}
 q&=\beta J+1,& C&=q-F-\alpha,\\
 X&=wq,&Y&=sq^3,&E&=XY,\\
 k&=\eta+\zeta,&c&=kY+\eta,&a&=E+Y,\\
 H&=4a+3,&\Delta&=a^2+H=(a+2)^2-1,\\
 \ell&=EkY,&A&=1+\ell(\ell+k),&
 G&=1+\Delta c^2.
\end{align*}
Set
\begin{align}
 R&=\bigl((q-1)(q-F)+\alpha\bigr)(q^2-1)
                  +(M_C+qM_F)J,\label{eq:R}\\
 N_k&=k-hE-R,\label{eq:Nk}\\
 N_t&=(K_0+w)C+q-F-z(q-1).\label{eq:Nt}
\end{align}
The remaining abbreviations are
\begin{align*}
 D&=1+\Delta i^2c^4,&Q&=(\Delta ic^2)^2,&d&=Q-1,\\
 v&=cT,&b&=c+RD,\\
 e&=Q(Dv^2+b^2)-1,&B_{\rm aux}&=2Qvb,\\
 W&=v^2+D+2b,&L&=1+HW,&S&=X+ac+Hr.
\end{align*}
The compiler base \(B\) and the auxiliary cross coefficient
\(B_{\rm aux}\) are different quantities.

\subsection{Quadratic norms as finite polynomial notation}

For a polynomial in a formal letter \(u\), reduce it modulo \(u^2-g\)
to \(p_0+p_1u\), and define
\[
 \Norm_{g,u}(p)=p_0^2-gp_1^2.
\]
All coefficients may themselves be polynomials. This operation is a
finite combination of integer additions and multiplications. It
introduces no quantified variable.

With temporary formal letters \(u,t\), put
\begin{equation}\label{eq:Fformal}
 \mathcal F(u,t)=
 \bigl[d(S-Lu-HLt)^2-H^2e\bigr]^2-H^4B_{\rm aux}^2D,
\end{equation}
and define the polynomial
\begin{equation}\label{eq:Nnorm}
 N=\Norm_{G,u}\bigl(\Norm_{A,t}(\mathcal F(u,t))\bigr).
\end{equation}
Thus \(u,t\) in \eqref{eq:Fformal}--\eqref{eq:Nnorm} are instructions
for constructing a polynomial, not two more unknowns.

The recommended single equation is
\begin{equation}\label{eq:product}
 \boxed{\quad P_{12}=N_kN_t(1+N)-1=0.\quad}
\end{equation}
Two useful alternatives are
\begin{equation}\label{eq:system}
 \boxed{\quad N_k=1,\qquad N_t=1,\qquad N=0\quad}
\end{equation}
and
\begin{equation}\label{eq:sos}
 Q_{12}=(N_k-1)^2+(N_t-1)^2+N^2=0.
\end{equation}
Every exponent in these definitions is fixed. The delivered JSON
circuits express them using only binary \(+,-,\times\).

\subsection{An explicit pair evaluator for the norm}

For readers wishing to avoid resultants or formal radicals altogether,
work with pairs of integer polynomials and the multiplication
\[
 (x_0,x_1)\circ(y_0,y_1)
   =(x_0y_0+Gx_1y_1,\ x_0y_1+x_1y_0).
\]
Addition, subtraction and scalar multiplication are coordinatewise.
This is multiplication modulo \(u^2-G\).

Put \(a_0=H^2L^2A\), \(Z=H^4B_{\rm aux}^2D\), and
\begin{align*}
 \mathcal V&=(S^2+L^2G,\ 2SL),\\
 \mathcal C&=d\bigl(\mathcal V+(a_0,0)\bigr)-(H^2e,0),\\
 \mathcal K&=4d^2a_0\mathcal V,\\
 (\mathcal N_0,\mathcal N_1)
   &=\bigl(\mathcal C^{\circ2}-\mathcal K-(Z,0)\bigr)^{\circ2}
                                    -4Z\mathcal K.
\end{align*}
Then
\begin{equation}\label{eq:paireval}
 N=\mathcal N_0^2-G\mathcal N_1^2.
\end{equation}
This is a closed polynomial recipe in the twelve supplied integers.
The identity behind the compact evaluator is
\[
 (C_0^2+K_1-Z)^2-4C_0^2K_1
       =(C_0^2-K_1-Z)^2-4K_1Z.
\]
It saves six binary gates in the pair implementation. The use of
\(S+Lu\) in \(\mathcal V\) instead of \(S-Lu\) is harmless because the
last norm includes both signs of \(u\).

\section{The rational auxiliary-Pell lemma}\label{sec:pell}

\begin{lemma}[Rational Pell ordinate is integral]\label{lem:rational}
Let \(m\ge2\) and \(V\) be integers. If \(y\in\Q\) satisfies
\[
 m^2V^2-(m^2-1)y^2=1,
\]
then \(y\in\Z\).
\end{lemma}

\begin{proof}
Write \(y=a_1/b_1\) in lowest terms, with \(b_1>0\).
Clearing denominators shows \(b_1^2\mid m^2-1\). Put
\[
 D_0=(m^2-1)/b_1^2.
\]
It is a positive nonsquare integer: if it were a square, then
\(m^2-1\) would be a square. We have two positive Pell solutions
\[
 m^2-D_0b_1^2=1,\qquad
 (m|V|)^2-D_0|a_1|^2=1.
\]
Neither ordinate is zero. Let
\(\epsilon=x_1+y_1\sqrt{D_0}>1\) be the least positive integral
norm-one Pell unit. Write
\[
 \epsilon^n=\chi_n+\psi_n\sqrt{D_0}.
\]
Elementary Pell classification gives indices \(j,l\ge1\) with
\[
 m=\chi_j,\quad b_1=\psi_j,\quad
 m|V|=\chi_l,\quad |a_1|=\psi_l.
\]

We need the precise first-coordinate divisibility statement
\begin{equation}\label{eq:pell-divis}
 \chi_j\mid\chi_l
 \quad\Longleftrightarrow\quad
 l\equiv j\pmod{2j}.
\end{equation}
Indeed, \(\chi_{2j}=2\chi_j^2-1\) and
\(\psi_{2j}=2\chi_j\psi_j\) imply
\(\chi_{n+2j}\equiv-\chi_n\pmod{\chi_j}\).
The subtraction identity similarly gives
\(\chi_{2j-n}\equiv-\chi_n\pmod{\chi_j}\).
Reduce \(l\) modulo \(2j\), then reflect into \(0\le n\le j\).
For \(n<j\), strict growth gives \(0<\chi_n<\chi_j\), so its signed
residue cannot vanish. The only possibility is \(n=j\), proving
\eqref{eq:pell-divis}.

Thus \(l\) is an odd multiple of \(j\). Pell multiplication gives
\(\psi_j\mid\psi_l\), hence \(b_1\mid a_1\). Coprimality forces
\(b_1=1\).
\end{proof}

The lemma is the specific integrality input needed below. In general,
\(dy^2\in\Z\) does not imply \(y\in\Z\), and neither does
\(Hy\in\Z\). The special relation between the two auxiliary Pell
coefficients is doing real work.

\section{Soundness of the norm reduction}\label{sec:soundness}

\subsection{The inherited arithmetic interface}

We state explicitly the interface imported from the source packages.
For an authentic coefficient tuple, the normalized blank compiler has
positive coordinates \(J,F,\alpha,z,i,s,w,\eta,\zeta\) and permits
signed \(h,T,\gamma,f,\tau,y\) once its signed-domain proof is applied.
Its four Pell conditions are
\begin{align}
 \tau^2-\ell(\ell+k)&=1,\label{eq:first}\\
 u^2-\Delta c^2&=1,\qquad u=X+ac+\gamma H,\label{eq:mainnorm}\\
 f^2-\Delta i^2c^4&=1,\label{eq:strong}\\
 QV^2-(Q-1)y^2&=1,\quad
 V=c(Tf-1)-Rf^2.\label{eq:aux}
\end{align}
Together with \(N_k=N_t=1\), these equations encode blank K215
halting. The same soundness argument allows a common unit sign
\(N_k=N_t\in\{1,-1\}\) initially, and proves that it is \(+1\).

The normalized strong equation maps to the ordinary strong interface
by \(i_{\rm ordinary}=\Delta i\). Its coefficient becomes
\[
 (i_{\rm ordinary}c^2)^2
   =\Delta(f^2-1)=Q.
\]
This is an integral substitution. We never assume that an arbitrary
ordinary multiplier can be divided by \(\Delta\).

The relevant full signed proofs are package 12's
\texttt{ORDINARY\_BLANK\_PROOF.md} and package 14's
\texttt{signed\_domain\_proof.md}; package 09 states the weak-sign
version used by its own product \cite{p09,p12,p14}. Copies and pinned paths are in the
archive. These are substantive inherited mathematical results,
not conclusions of the finite tests performed here.

\subsection{Bootstrap before reconstructing the roots}

Initially suppose \(N=0\) and \(N_k=N_t=\varepsilon\), where
\(\varepsilon\in\{1,-1\}\). This covers the system, and will cover
the product after the parity argument in Section~\ref{sec:parity}.

The transport factor can be rewritten
\[
 N_t=(K_0+w)C+1-F-(z-1)(q-1).
\]
Since \(F,z\ge1\), if \(C\le-1\), then
\(N_t\le-(K_0+w)\le-2\), a contradiction. Thus \(C\ge0\).
For \(N_t=1\), the stronger equality gives \(C>0\); we retain
\(C\ge0\) so the proof also covers the negative provisional sign.

The authentic mask bounds now give
\[
 q\ge16,\qquad 0<R<q^4-2,\qquad E\ge q^4>R+2.
\]
Consequently \(A,G,D,Q>1\), \(d>0\), and \(b,v,W,L,H>0\).
These bounds precede the recovery of the main root, so there is no
circular assumption about its sign or index.

\subsection{One coherent branch of the norm}

Over an algebraic closure, the evaluated polynomial \(N\) is exactly
\begin{equation}\label{eq:signproduct}
 \prod_{\substack{u=\pm\sqrt G\\t=\pm\sqrt A\\f=\pm\sqrt D}}
 \bigl[d(S-Lu-HLt)^2-H^2e+H^2B_{\rm aux}f\bigr].
\end{equation}
If it is zero, choose one vanishing factor and put
\[
 y=(S-Lu-HLt)/H.
\]
The chosen roots satisfy
\begin{align}
 u^2&=G,\qquad t^2=A,\qquad f^2=D,\label{eq:roots}\\
 dy^2&=e-B_{\rm aux}f=Q(vf-b)^2-1,\label{eq:yrelation}\\
 S&=Lu+HLt+Hy.\label{eq:separated}
\end{align}
The signs are chosen from one actual zero factor. We do not assume
that all independent formal sign choices occur as field automorphisms.

\subsection{Uniform bounds for every conjugate}

Take a finite normal extension of \(\Q\) containing this selected
branch. Under every embedding, \(f\) goes to a square root of \(D\).
For every conjugate \(y'\), including a complex one,
\begin{align*}
 |y'|^2
 &\le \frac{e+B_{\rm aux}\sqrt D}{d}\\
 &=\frac{Q(v\sqrt D+b)^2-1}{Q-1}
 <2(v\sqrt D+b)^2.
\end{align*}
Since \(v^2+D\ge2v\sqrt D\),
\begin{equation}\label{eq:ybound}
 |y'|<W=v^2+D+2b.
\end{equation}
In particular,
\begin{equation}\label{eq:firstsep}
 L\sqrt G>HW.
\end{equation}

The other separation estimate is
\begin{equation}\label{eq:secondsep}
 HL\sqrt A>L\sqrt G+HW.
\end{equation}
Here is a direct proof. We have \(\sqrt A>\ell\) and
\(\sqrt G<(a+2)c\), because \(c>1\) and
\(\Delta=(a+2)^2-1\). Also
\(c<k(Y+1)\le2kY\). Therefore
\begin{align*}
 H\sqrt A-\sqrt G
 &>H\ell-(a+2)c\\
 &>kY\bigl(HE-2(a+2)\bigr)\\
 &\ge kY(H-2a-4)=kY(2a-1)>1.
\end{align*}
Multiplying by \(L=1+HW>HW\) proves \eqref{eq:secondsep}.
Thus one weight \(L\), used with the factors \(1,H\), separates
both roots. A second large independent weight is unnecessary.

\subsection{Rationality, then integrality}

If \(t\) were irrational, an automorphism of the normal extension
would send \(t\) to \(-t\). Subtract its image of
\eqref{eq:separated} from the original. The changed \(t\)-term has
magnitude \(2HL\sqrt A\), whereas the other terms have total
magnitude at most \(2L\sqrt G+2HW\), contradicting
\eqref{eq:secondsep}. Thus \(t\in\Q\), and \(t^2=A\in\Z\)
then implies \(t\in\Z\).

If \(u\) were irrational, an automorphism would send it to \(-u\)
and fix the now rational \(t\). The same subtraction would give
\(2L\sqrt G\le2HW\), contrary to \eqref{eq:firstsep}.
Hence \(u\in\Z\). Equation~\eqref{eq:separated} now makes
\(y\) rational. Since \(B_{\rm aux}>0\),
\[
 f=\frac{e-dy^2}{B_{\rm aux}}\in\Q,
\]
and \(f^2=D\in\Z\) gives \(f\in\Z\).

Now \(V=vf-b\) is an integer and
\(Q=(\Delta ic^2)^2\). Lemma~\ref{lem:rational}, applied to
\eqref{eq:yrelation}, proves \(y\in\Z\).
This final step is the point at which rational recovery alone would
otherwise have been insufficient.

\subsection{Restore the main congruence}

Because \(L=1+HW\), equation~\eqref{eq:separated} gives
\[
 u-X-ac=H\bigl(r-Wu-Lt-y\bigr).
\]
Thus
\[
 \gamma=r-Wu-Lt-y\in\Z
\]
restores the main-root coordinate exactly.
Furthermore \(f^2=D\) implies
\[
 vf-b=c(Tf-1)-Rf^2.
\]
The recovered \(t,u,f,y,\gamma\) therefore satisfy all four
conditions \eqref{eq:first}--\eqref{eq:aux}, with \(\tau=t\).
Their signs need not be guessed. The signed ordinary interface,
after the integral multiplier conversion stated above, applies.
It recovers the main root's positive sign, forces the provisional
index/transport signs to \(+1\), and decodes a blank halting computation.
This establishes soundness.

\section{Completeness and the single-product finalizer}

\subsection{Every canonical halting witness gives a new witness}

If K215 halts on the blank tape, the inherited inputless compiler
produces a canonical positive normalized witness. Write its positive
main root as \(u=X+ac+\gamma H\) and its first root as \(\tau\).
Keep the eleven coordinates occurring in the new definitions and set
\begin{equation}\label{eq:complete-r}
 r=\gamma+Wu+L\tau+y.
\end{equation}
This is positive and satisfies
\[
 S=Lu+HL\tau+Hy.
\]
One factor of \eqref{eq:signproduct} is zero. Hence \(N=0\),
while \(N_k=N_t=1\) already hold. This proves completeness for
the system, the product and the sum of squares.

The witness transformation deletes \(f,\tau,y,\gamma\) and adds \(r\):
fifteen supplied coordinates become twelve. It makes no claim that
the new \(r\) is numerically small.

\subsection{Why \texorpdfstring{\(1+N\)}{1+N} is a valid unit guard}\label{sec:parity}

The pair evaluator writes the inner norm as
\(\mathcal N_0+\mathcal N_1u\).
Its odd coefficient \(\mathcal N_1\) is divisible by \(8\) as an
integer polynomial. One way to see this is to put \(x=S+Lu\).
The odd coefficient of \(x^2\) is even, that of \(\mathcal C^2\)
is divisible by \(4\), and that of \(\mathcal K\) is divisible
by \(8\). Thus the odd coefficient of
\[
 (\mathcal C^2+\mathcal K-Z)^2-4\mathcal C^2\mathcal K
\]
is divisible by \(8\). This expression equals the optimized
pair expression in Section~\ref{sec:formula}.
Consequently
\[
 N=\mathcal N_0^2-G\mathcal N_1^2
      \equiv \mathcal N_0^2\pmod{64}.
\]
In particular \(1+N\) cannot be \(-1\) modulo \(4\).

At a zero of \(P_{12}\), each of the three integer factors
\(N_k,N_t,1+N\) is a unit. The third must be \(1\), so \(N=0\)
and \(N_k=N_t\in\{1,-1\}\). The soundness proof was deliberately
established for exactly this provisional sign pattern. It forces
both signs positive.

\begin{theorem}\label{thm:positive}
For either the new complete coefficient tuple or the inherited T4
tuple, the three-equation system \eqref{eq:system},
the equation \(P_{12}=0\), and the equation \(Q_{12}=0\) have
positive integer solutions exactly when blank K215 halts.
Their no-solution statements are therefore each provably equivalent
in ZFC to \(\Con(\mathrm{ZFC})\).
\end{theorem}

\section{Exact degree and circuit size}\label{sec:degree}

Every supplied unknown has degree one; fixed coefficients have
degree zero. The useful degree ledger is
\begin{center}
\begin{tabular}{lrrrrrrrr}
\toprule
Expression & \(c\) & \(a,H\) & \(\Delta\) & \(R\) & \(\ell\) & \(A,G\) & \(Q\) & \(D\)\\
Degree & 5 & 6 & 12 & 4 & 11 & 22 & 46 & 34\\
\midrule
Expression & \(v\) & \(b,W\) & \(L\) & \(S\) & \(a_0\) & \(\mathcal V_0\) & \(\mathcal V_1\) & \(\mathcal C_0\)\\
Degree & 6 & 38 & 44 & 11 & 122 & 110 & 55 & 168\\
\bottomrule
\end{tabular}
\end{center}
The leading term of \(\mathcal C_0\) comes uniquely from \(da_0\).
The even coefficient \(\mathcal N_0\) has degree \(672\), with
leading term \((\mathcal C_0)_{\topdeg}^4\); the odd coefficient
has degree at most \(605\). Thus the term
\(\mathcal N_0^2\) has degree \(1344\), whereas
\(G\mathcal N_1^2\) has degree at most \(1232\).
There is no top-degree cancellation.

More explicitly, set \(q_0=\beta J\), \(k_0=\eta+\zeta\). Then
\begin{equation}\label{eq:Ntop}
 N_{\topdeg}
  =2^{80}i^{48}k_0^{112}w^{112}s^{224}
             q_0^{832}(q_0-F)^{16}.
\end{equation}
It is nonzero for every \(\beta>0\).
The leading forms of the outer factors are
\[
 (N_k)_{\topdeg}=-hwsq_0^4,\qquad
 (N_t)_{\topdeg}=w(q_0-F-\alpha)-zq_0.
\]
Therefore \(\deg P_{12}=1344+7+2=1353\).
For example, the coefficient of
\[
 J^{853}hz\,i^{48}\eta^{112}w^{113}s^{225}
\]
is \(2^{80}\beta^{853}\ne0\). The sum of squares has exact
degree \(2688\).

The literal circuit counts are
\begin{center}
\begin{tabular}{lrrr}
\toprule
Presentation & Multiplications & Additions/subtractions & Total\\
\midrule
Three equations & 71 & 41 & 112\\
Product polynomial & 73 & 43 & 116\\
Sum of squares & 74 & 45 & 119\\
\bottomrule
\end{tabular}
\end{center}
The generator checks that every gate, every variable and all four
coefficient ports are live. These counts concern the delivered
evaluation schedule; no optimal arithmetic-circuit claim is made.

\subsection{An optional lower-degree weight}

There is a second twelve-positive-variable polynomial that exchanges
seven extra operations for lower degree. Keep all other definitions,
but replace \(W\) by
\[
 W'=4T(ic^2)+2+2R(a+6)i^2c^3,\qquad L'=1+HW'.
\]
The norm construction and product finalizer are unchanged.
This gives a product of exact degree \(1177\) at \(123\) operations,
a three-equation system of maximum degree \(1168\) at \(119\)
operations, and a sum of squares of degree \(2336\) at \(126\)
operations.

The shorter degree comes from estimating conjugates relative to
\(\sqrt G\), instead of absolutely. Weak outer soundness and \(h>0\)
give \(R<c\). With \(I=ic^2\), use
\(\sqrt G>ac\), \(\sqrt D<(a+2)I\), and
\(\Delta/a=a+4+3/a<a+6\) to obtain
\[
 \frac{|y'|}{\sqrt G}
 <\frac{2TI(a+2)}a+\frac2a+\frac{2R}{ac}
                       +\frac{2R\Delta i^2c^3}{a}
 <W'.
\]
Also \(H\sqrt A>2\sqrt G\), since \(E\ge2\) and
\(HE\ge2H>4(a+2)\). These are sufficient for the same two
automorphism contradictions, because
\[
 H|y'|<L'\sqrt G,\qquad
 HL'\sqrt A>2L'\sqrt G>L'\sqrt G+H|y'|.
\]
The integrality lemma, main congruence and product guard are unchanged.
Now \(\deg W'=27\) and \(\deg L'=33\); the norm's leading form is
\[
 2^{80}i^{48}k_0^{96}w^{96}s^{192}q_0^{720}(q_0-F)^{16}.
\]
It has degree \(1168\), proving the product degree \(1177\).
This alternate weight is proved here only for positive supplied
coordinates. The \(26\)-integer extension in the next section uses
the main weight \(W=v^2+D+2b\).

\section{A polynomial in 26 unrestricted integers}\label{sec:integer}

\subsection{A larger admissible signed domain}

The same equations are sound with only the following seven coordinates
required to be positive:
\[
 J,F,\alpha,z,s,\eta,\zeta.
\]
Allow \(h,T,r\) to be arbitrary integers, assume \(i\ne0\), and
allow \(w\ge0\). This extension needs two boundary arguments; it is
not obtained just by changing the declared domain.

Transport still gives \(C\ge0\) and \(R>0\), so \(b>0\).
The conjugate bound remains valid with \(|v|\) in place of \(v\):
\[
 v^2+D+2b\ge2(|v|\sqrt D+b)>|y'|.
\]
If \(w>0\), the same separation proof recovers rational \(t,u,y\).
If \(w=0\), then \(E=\ell=0\), so \(A=1\) and \(t=\pm1\) is
already rational. The second stage of separation still proves that
\(u,y\) are rational.

When \(T\ne0\), recover \(f\in\Z\) as before and use
Lemma~\ref{lem:rational}. When \(T=0\), \(v=B_{\rm aux}=0\);
instead the root relation directly reads
\[
 (Q-1)y^2=Qb^2-1.
\]
The same lemma, with \(V=b\) and \(m=|\Delta ic^2|\), still
proves \(y\in\Z\). Thus the main congruence and integral
\(\gamma=r-Wu-Lt-y\) are available in all cases.

If \(w=0\), then \(X=0\), \(a=Y\), and \(u=ac+\gamma H\).
The main norm would give
\[
 1=u^2-\Delta c^2
   =H\bigl(2ac\gamma+H\gamma^2-c^2\bigr),
\]
impossible because \(H>1\). Therefore \(w>0\).

Now the authentic bounds give \(R<q^4\) and
\(\Delta>a^2\ge q^6\), hence \(\Delta>R\).
For \(T=0\), the integer \(b=c+RD\) satisfies \(1<b<Q\), because
\[
 Q-b=\Delta i^2c^4(\Delta-R)-c-R>0
\]
using \(c\ge2\). But then
\[
 y^2=b^2+\frac{b^2-1}{Q-1}
\]
lies strictly between \(b^2\) and \((b+1)^2\), impossible.
Thus \(T\ne0\), so \(f\) is recovered in every actual solution.

The inherited signed interface now applies with the positive ordinary
multiplier \(\Delta|i|\), and proves blank halting. The sign of \(h\)
is harmless: once \(E>R+2\), the relation
\(k=R+\varepsilon+hE>0\) forces \(h\ge0\), after which the
usual index argument applies.

\subsection{Completeness with odd \texorpdfstring{\(i\)}{i} and square \texorpdfstring{\(w\)}{w}}

The source package 14 proves a canonical witness can be chosen with
\(i\) odd and \(w\) a nonzero square \cite{p14}. We recall why the new coefficient
recipe preserves its prerequisites.
The spatial width and temporal height may both be sufficiently large
powers of five. The coefficient exponent \(d_{\rm cell}=5^{15}\)
is odd, so the exponent \(t_{\rm pow}\) in \(q=2^{t_{\rm pow}}\)
is odd. The canonical tuned rank is \(R\equiv3\pmod4\), hence odd,
and
\[
 w=2^{R-t_{\rm pow}}
\]
is a square.

For the normalized strong witness, let
\(c=\psi_{a+2}(R)\), which is odd, and choose the auxiliary
index \(m=cR\). Then \(c^2\mid\psi_{a+2}(cR)\) by the binomial
expansion of the \(c\)-th power of the Pell unit at index \(R\).
Thus
\[
 i=\psi_{a+2}(cR)/c^2
\]
is a positive odd integer. The same odd-index quotient identities
give integral positive auxiliary witnesses, as proved in package 14.
The new language permits precisely the padding and five-adic tuning
used in that proof. Finally define \(r\) by \eqref{eq:complete-r}.

\subsection{The explicit substitution}

Use the positive surjection
\[
 \operatorname{Pos}(x,y,t)=1+x^2+y^2+t(t+1).
\]
It is positive for all integers and represents every positive integer.
For \(n\ge1\), the three-square theorem represents \(4n-3\).
Modulo four exactly one square is odd and two are even; writing
them as \((2t+1)^2,(2x)^2,(2y)^2\) gives the formula above.

Replace each of the seven positive ports by its own
\(\operatorname{Pos}\), and set
\[
 i=2j+1,\qquad w=u_0^2,
\]
leaving \(h,T,r\) unrestricted. The resulting polynomial
\(P_{26}^{\Z}\) has
\[
 7\cdot3+5=26
\]
unrestricted integer variables. The soundness and completeness
arguments just given show
\[
 \mathrm{ZFC}\vdash
 \left(\forall\mathbf z\in\Z^{26}\ P_{26}^{\Z}(\mathbf z)\ne0\right)
 \longleftrightarrow\Con(\mathrm{ZFC}).
\]
The seven positive substitutions cost \(49\) gates; the odd multiplier
and square scale cost three more. Hence the product uses \(168\)
operations. The corresponding weighted degrees are
\[
 \deg N=2640,\qquad \deg N_k=13,\qquad \deg N_t=4,
\]
giving exact product degree \(2657\).
The degree proof is uniform in \(\beta>0\), and the delivered audit
records a nonzero leading monomial.

\section{The smaller complete coefficient tuple}\label{sec:coefficients}

\subsection{A stricter finite language}

The machine remains the literal \(215\)-state K215, with \(430\)
transitions and a separate halt state. Its dense table has SHA-256
\begin{center}\scriptsize
\texttt{fdbcef0825886035c3091c8537a69535bc96f9e265f0ae45863543d40ae70090}.
\end{center}
Halt is completed by a stationary transition only for padding an
already halted computation; the original transitions entering halt
are unchanged.

Retain the \(439\) cell symbols from T4 \cite{p11}: headless bits \(0,1\);
headed code \(2(q+1)+b\), \(0\le q\le215\), \(b\in\{0,1\}\);
initialization symbols \(L_0=434,Q_0=435,R_0=436\);
vertical boundary \(V_0=437\); horizontal boundary \(H_0=438\).
A tile is \((l,c,r,n)\), representing left, center, right, and next.

For every tile other than the special Forbidden tile, impose the
old T4 predicate and three additional guards:
\begin{enumerate}
 \item Both pairs \((l,c)\), \((c,r)\) obey the symmetric row rule.
 Away from a vertical boundary, symbols have the same kind:
 ordinary, initial phase, or horizontal boundary. Initial-phase
 pairs are \(L_0L_0,L_0Q_0,Q_0R_0,R_0R_0\). A pair containing
 \(V_0\) is allowed unless both entries are \(V_0\).
 \item If \(V_0\) occurs among \(l,c,r\), every ordinary symbol
 in that triple is blank and headless.
 \item If \(n=H_0\), every head among \(l,c,r\) is in halt state \(215\).
\end{enumerate}
Retain Start \(=(434,435,436,1)\) first and Forbidden
\(=(438,0,0,438)\) second. Forbidden remains an old-valid local
tile, but its selector is prohibited by the mask. Append all other
strict tiles in lexicographic order.

There are \(5279\) strict tiles and \(5280\) including Forbidden.
An exhaustive enumeration and a separate center-class census agree:
\begin{center}
\begin{tabular}{lrrrrrrrr}
\toprule
Center & Headless & Working head & Halt head & \(L_0\)&\(Q_0\)&\(R_0\)&\(V_0\)&\(H_0\)\\
Count &3495&1720&16&6&1&6&24&12\\
\bottomrule
\end{tabular}
\end{center}

Every new tile is an old tile, so soundness is inherited by taking a
subset. For completeness, pad a finite blank halting computation
with at least three never-visited blank cells at each vertical wall,
and close its last row after halting. The new guards hold throughout,
including helical seams. Arbitrarily large independent padding in
width and height is still possible, including powers of five.

\subsection{All fixed scalars and finite sums}

In this subsection use \(A_{\rm cl}\) for the clause radix, to avoid
confusion with the first Pell radicand \(A\).
Set
\[
 a_{\rm cell}=439,\quad k_{\rm tile}=5280,\quad
 A_{\rm cl}=8192,\quad \rho=2^{78125},\quad
 B=\rho^{390625}=2^{5^{15}}.
\]
For \(0\le u\le438\), define
\[
 e_l(u)=7036+u,\quad e_c(u)=6597+u,\quad
 e_r(u)=6158+u,\quad e_n(u)=5280+u.
\]
The support \(\mathcal E\) is the union of selectors \(0,\ldots,5279\),
these four copy intervals, dummies \(7475,7476\), and the four
ordered anchors
\[
 7480,\ 22440,\ 14960,\ 44880.
\]
There are \(7042\) support positions.
For tile \(w_j=(l_j,c_j,r_j,n_j)\), put
\begin{align*}
 v_j={}&1+A_{\rm cl}^{1+l_j}+A_{\rm cl}^{440+c_j}
          +A_{\rm cl}^{879+r_j}+A_{\rm cl}^{1318+n_j}\\
       &+\sum_{t=1}^4 A_{\rm cl}^{1756+t}.
\end{align*}
At the four copy intervals use respectively
\[
 A_{\rm cl}^{1+u},\quad A_{\rm cl}^{440+u},\quad
 A_{\rm cl}^{879+u},\quad A_{\rm cl}^{1318+u}.
\]
At the ordered anchors use \(A_{\rm cl}^{1756+t}\), \(t=1,\ldots,4\);
at the dummies use zero. This specifies \(v_e\) at every \(e\in\mathcal E\).
Let
\[
 \mu=A_{\rm cl}-2+\sum_{j=1}^{2849}A_{\rm cl}^j.
\]
The complete constants are
\begin{align*}
 D_C={}&\rho^{1317}+\rho^{69077}+\rho^{7480}+\rho^{22440}\\
       &+\sum_{e\in\mathcal E}v_e
              \bigl(\rho^{158838-e}+\rho^{248600-e}\bigr),\\
 D_R={}&\rho^{68638},\\
 M_F^0={}&\mu(\rho^{158838}+\rho^{248600})
                   +\sum_{u=0}^{438}\rho^{6597+u}\\
       &+\sum_{u=0}^{438}
          \bigl(\rho^{68638+7036+u}+\rho^{68638+6597+u}\bigr)
          +\rho^{14960}+\rho^{44880},\\
 M_C={}&B-1-\sum_{\substack{e\in\mathcal E\\e\ne1}}\rho^e
                                       -2\rho^{7475},\\
 M_F^{\rm nat}={}&M_F^0+4.
\end{align*}
Finally,
\begin{equation}\label{eq:tuple}
 \boxed{\quad
 (\beta,K_0,M_C,M_F)
   =(B-1,\ D_C+BD_R,\ M_C,\ M_F^{\rm nat}+B-1).
 \quad}
\end{equation}
These are four particular integers, defined by bounded finite
enumerations and sums. Their definition never asks whether K215
halts. The JSON recipe and literal alphabet fix every ordering choice.

\subsection{Why the modified packing still works}

The copy intervals are packed more closely than in T4. The new anchor
set is \(\{M,2M,3M,6M\}\), with \(M=7480\).
The anchor tests map \(M\) to \(2M\) and \(3M\) to \(6M\),
using shifts \(M,3M\). The difference \(4M\) has exactly one
ordered realization in the whole support, namely \((2M,6M)\).
All nonanchor positions are below \(M\), and the cell circumference
is more than twice the support diameter. Therefore a whole-radix
rotation preserving both target anchors has zero cell displacement.

The scalar geometry is
\[
 E_{\max}=44880,\quad H_c=68638,\quad
 T_1=158838,\quad T_2=248600,\quad L_{\rm cell}=390625.
\]
The two clause bands have gap \(89762>2E_{\max}\), congruent to
\(2\pmod4\), and
\[
 T_2+2E_{\max}+2=338362<L_{\rm cell}.
\]
Thus at least one clause band is clean under every support rotation.
The inherited no-carry and odd-anchor argument first forces a binary
rotation to be a radix rotation; the unique anchor difference then
forces whole-cell synchronization.

The verifier checks all \(1319=3\cdot439+2\) copy and anchor
convolution targets by exact finite set identities. It also checks
clause-band isolation, the full support population, and the original
conservative carry bound. The clause target has population
\(\pc(\mu)=2861\), so
\[
 \pc(M_F^{\rm nat})=7042,\qquad
 \pc(M_C)=5^{15}-7042.
\]
The radix exponent \(78125\) exceeds both the \(37053\)-bit
conservative threshold and the direct coefficient-mass threshold.

Both \(78125=5^7\) and \(390625=5^8\) are powers of five.
Because \(M\) is divisible by four and \(T_2-T_1\equiv2\pmod4\),
the actual coefficient residue remains
\[
 2D_C-D_R\equiv3\pmod5.
\]
This is checked by summing the actual finite clause coefficients,
as well as by the band-cancellation argument.
Consequently the original dummy-bit subset adjustment still tunes
the actual packed rank to the required congruence modulo
\(d_{\rm cell}N_{\rm cell}\), with \(d_{\rm cell}=5^{15}\).
The second dummy need not be used.

The full adapter proof is in
\texttt{proofs/REDUCED\_T4\_PROOF.md}. It identifies which parts are
inherited compiler lemmas and which are new finite support identities.
This preserves the inputless compiler interface used above.

\section{Reproducibility, evidence, and limits}\label{sec:repro}

\subsection{What is supplied}

The archive contains this \LaTeX\ source and its compiled PDF,
the nine JSON arithmetic circuits, the two standalone circuit generators,
the complete \(5280\)-tile alphabet, the coefficient recipe,
the literal K215 transition table and JSON, the finite compiler
checker, independent proof/audit notes, selected pinned upstream
proofs, source hashes, and a single reproduction entry point.
The inherited 67-operation circuit is also included as
\texttt{upstream/PRODUCT67.json}; the coefficient binding records its
port aliases for the new tuple.
The arithmetic and finite checks require only Python's standard library.

From the extracted package directory:
\begin{verbatim}
python3 verify.py
latexmk -pdf -interaction=nonstopmode -halt-on-error report.tex
\end{verbatim}
The verification driver regenerates and checks the circuits and finite
alphabet. It does not expand the final four huge integers.

\subsection{Checks and their interpretation}

The algebra checks compare the norm circuit with a separately
implemented product of all eight conjugates in an eight-dimensional
multiquadratic quotient algebra. They include signed off-solution
assignments, parity identities, all-port liveness, exact operation
counts, integer-substitution identities, and leading-form certificates.
The uniform symbolic leading forms, rather than numerical sampling,
prove the stated exact degrees.

The coefficient checker binds all \(430\) machine transitions to the
literal table hash, examines all \(46980\) horizontal triples surviving
the one-head filter, and tests all \(439\) successors for the \(13465\)
triples surviving the row filter. This is \(5,911,135\) explicitly
tested quadruples. Every skipped triple violates a stated guard.
It checks the class count, every convolution target, the support and
anchor conditions, mask populations, carry estimates, and actual
coefficient residue modulo five. It also tests \(1143\) helical
windows from two deliberately halting control machines.

These checks support concrete finite and algebraic claims.
They do not establish the unbounded compiler theorem by testing,
provide an accepting K215 computation, or prove ZFC consistent.

\subsection{Formalization boundary}

The reviewed K215 verification records report the exact dense-machine
endpoint in Lean and Rocq \cite{endpoint}. We fetched and inspected the endpoint and
its recorded scope; we did not rerun those large builds.
No new end-to-end kernel proof of the norm reduction or coefficient
adapter is claimed. The delivered proof is an ordinary mathematical
proof composed with explicitly named inherited lemmas.

There is no global optimality claim. In particular the repository
already discusses an abstract nine-unknown existence theorem derived
from classical universal representations \cite{repo}. This report supplies a
smaller explicit instance in the concrete compiler family; it does
not improve the universal variable frontier. It also does not improve
the family's \(67\)-operation minimum or its low-degree systems.

Finally, \eqref{eq:main-con} is an equivalence, not a resolution of
\(\Con(\mathrm{ZFC})\). Under the usual incompleteness hypotheses,
consistency gives unprovability of the no-solution statement in ZFC.
Two-sided undecidability requires an additional soundness hypothesis;
we do not infer it from consistency alone.

\section*{Sources and attribution}
\addcontentsline{toc}{section}{Sources and attribution}

All repository sources below are pinned to \pin. The accompanying
\texttt{SOURCE\_MANIFEST.json} gives exact paths and hashes of included
source copies. The upstream ProveIt and machine notices are preserved
in the archive.

\begin{thebibliography}{99}
\raggedright
\bibitem{repo}
Vladimir Reshetnikov et al.,
\emph{ProveIt: Con(ZFC) in Diophantine form}, source snapshot above.
\url{https://github.com/VladimirReshetnikov/ProveIt}.
The reconciliation README in the named report directory identifies
packages 09--15 and their stated metrics.

\bibitem{p09}
Package 09, \emph{Smaller Diophantine descriptions of ZFC consistency}.
Files \texttt{con\_zfc\_reductions.tex},
\texttt{proofs/ROOT\_AMALGAMATION\_PROOF.md}, and the accompanying circuits.
Weighted three-root amalgamation and weak sign recovery.

\bibitem{p11}
Package 11, \emph{Combining explicit Diophantine encodings of ZFC consistency}.
Files \texttt{bridge/BRIDGE\_AUDIT.md},
\texttt{bridge/T4\_CONSTANTS.json}, and \texttt{bridge/check\_t4.py}.
The original T4 coefficient compiler.

\bibitem{p12}
Package 12, \emph{Smaller Diophantine forms of ZFC consistency}.
\texttt{proofs/ORDINARY\_BLANK\_PROOF.md} and
\texttt{proofs/ROOT\_ELIMINATION\_PROOF.md}.
Ordinary strong Pell interface and signed main-root soundness.

\bibitem{p14}
Package 14, \emph{Con(ZFC) synthesis}.
\texttt{proofs/signed\_domain\_proof.md},
\texttt{proofs/independent\_signed\_review.md}.
The seven-positive-port extension, odd multiplier, and square scale.

\bibitem{blank}
Retained wave-1 package 06 proofs, bundled in the current packages:
\texttt{blank\_helical\_proof.md},
\texttt{fixed\_blank\_67\_proof.md}, and
\texttt{pell\_lemma\_audit.md}.
Blank-input computation tableaux, actual-rank tuning, and Pell lemmas.

\bibitem{endpoint}
ProveIt, \texttt{Computability/BusyBeaver/Lean/BusyBeaver/ZFC/}\\
\texttt{KEndpoint215.lean}, and
\texttt{Computability/BusyBeaver/Research/ZFC/}\\
\texttt{k225-formal/verification/k215/README.md}.
Exact-machine calibration and formal verification records.

\bibitem{new}
This package: \texttt{proofs/NORM12\_PROOF.md},
\texttt{proofs/NORM12\_INDEPENDENT\_AUDIT.md},
\texttt{proofs/INTEGER26\_PROOF.md},
\texttt{proofs/REDUCED\_T4\_PROOF.md}, and the generated receipts.
New mathematical derivations and finite verification evidence.
\end{thebibliography}
\end{document}
